%% parts/p05/provability-logic-gl.tex --- Part 5
\chapter{Provability Logic}\label{sec:p05-provability-logic}

This chapter uses the explicitly selected \(\PA\) proof formula and the
derivability and L\"ob results developed above. The modal and arithmetical
completeness results have their separate source treatments in
\cite{Segerberg1971,Solovay1976,Boolos1993}.

\begin{hbdefinition}{def:p05-gl}{The provability logic \(\mathsf{GL}\)}
The language is propositional modal logic with \(\Box\).  The Hilbert calculus
\(\mathsf{GL}\) contains classical propositional logic, the normal axiom
\(\Box(A\to B)\to(\Box A\to\Box B)\), theorem-only necessitation---from
\(\mathsf{GL}\vdash A\) infer \(\mathsf{GL}\vdash\Box A\)---and L\"ob's axiom
\[
 \Box(\Box A\to A)\to\Box A.
\]
The name is necessary to fix this exact calculus rather than any logic with an
informal provability reading.  It stands as a modal consequence system,
differs from \(\mathsf K\) by L\"ob's axiom, and is used for both relational
and arithmetical completeness below \cite{Segerberg1971,Boolos1993}.  The
complete local rule presentation is Sheet~\ref{sheet:hilbert-gl};
necessitation is not a rule from an open assumption.
\index{GL@\(\mathsf{GL}\)}\index{Lob axiom@L\"ob axiom}
\end{hbdefinition}

\begin{hbdefinition}{def:p05-gl-kripke}{Kripke semantics for \(\mathsf{GL}\)}
A Kripke model is a triple \((W,R,V)\), where \(W\ne\varnothing\),
\(R\subseteq W^2\), and \(V\) assigns to each propositional letter a subset
of \(W\).  Boolean satisfaction is classical, and the modal clause is
\[
 (W,R,V),w\entails\Box A
 \quad\Longleftrightarrow\quad
 (\forall v\in W)(wRv\to(W,R,V),v\entails A).
\]
The chosen orientation therefore reads \(wRv\) as “\(v\) is accessible from
\(w\).”  A frame is \emph{conversely well-founded} when it has no infinite
chain \(w_0Rw_1Rw_2R\cdots\).  These clauses are needed to make the direction
and well-foundedness condition in the frame theorem unambiguous.  They define
a relational semantics, not formal derivability in arithmetic, and are used
only in the frame-completeness and finite-model arguments below.
\index{Kripke semantics!provability logic}
\end{hbdefinition}

\begin{hbdefinition}{def:p05-arithmetical-realization}{Arithmetical realization}
Fix \(\PA\) and the literal selected formula
\(\mathsf{Prv}^{19}_{\PA}(x)=\exists p\,\Gamma_{d_{19}}(p,x,\bar1)\)
from \cref{def:p05-selected-direct-pa-graph}. An arithmetical
realization \(^*\) assigns each propositional letter a closed
\(\Lang_A\)-sentence, commutes with Boolean connectives, and puts
\[
 (\Box A)^*=\mathsf{Prv}^{19}_{\PA}
   (\overline{\ulcorner A^*\urcorner}).
\]
The term is needed because modal validity and provability in one arithmetical
theory are different relations.  It stands as a translation, differs from a
Kripke valuation by replacing box with a coded predicate, and is used in
Solovay completeness.  No generic theory \(T\) is covered by this definition.
\index{arithmetical realization}\index{provability predicate}
\end{hbdefinition}

The Hilbert--Bernays--L\"ob derivability conditions are proved for this
same predicate in
\cref{thm:p05-direct-pa-internal-mp,thm:p05-direct-pa-d3}. Together with
L\"ob's theorem, they give arithmetical soundness of \(\mathsf{GL}\)
for this predicate.
This does not authorize
necessitation for arbitrary true but unproved sentences; the rule is applied
only to theorems of the selected formal theory.

\begin{hbtheorem}{thm:p05-gl-frames}{cited=Segerberg1971}{Finite-frame completeness of \(\mathsf{GL}\)}
\(\mathsf{GL}\) is sound and complete for the finite irreflexive transitive
Kripke frames under Definition~\ref{def:p05-gl-kripke}.  Equivalently, these
are finite strict partial orders in the displayed accessibility direction.
Finiteness already rules out infinite \(R\)-chains; in the unrestricted frame
formulation the corresponding additional condition is converse
well-foundedness.
\end{hbtheorem}

The finite-model property yields decidability.  It does not identify the modal
frame relation with formal derivability inside arithmetic; the relational
semantics is an independent completeness presentation.

\begin{hbtheorem}{thm:p05-modal-fixed-point}{cited=Sambin1976}{Modal fixed-point theorem}
If every occurrence of \(p\) in \(A(p)\) is within the scope of a box, then
there is a formula \(F\) not containing \(p\) such that
\(\mathsf{GL}\vdash F\leftrightarrow A(F)\).  The modalized-occurrence
hypothesis cannot be deleted.
\end{hbtheorem}

\begin{hbtheorem}{thm:p05-solovay-completeness}{cited=Solovay1976}{Solovay arithmetical completeness}
For the selected \(\mathsf{Prv}^{19}_{\PA}\) and exactly the realizations of
Definition~\ref{def:p05-arithmetical-realization},
\[
 \mathsf{GL}\vdash A
 \quad\Longleftrightarrow\quad
 \PA\vdash A^*
 \text{ for every arithmetical realization }{^*}.
\]
The ambient metatheory uses the standard model
\(\mathcal N\entails\PA\), and hence the required external consistency and
\(\Sigma^0_1\)-soundness of this specific theory.  The display is not a theorem
about an arbitrary consistent or recursively axiomatized \(T\); changing the
base theory, numbering, or proof predicate requires a separately stated
generalization.
\end{hbtheorem}

\begin{proofboundary}{def:p05-solovay-boundary}{Solovay completeness proof boundary}
The forward implication follows from the derivability conditions and L\"ob's
theorem.  For the converse, Solovay assigns arithmetical sentences to nodes of
a finite counterframe and constructs a self-referential limit process whose
possible terminal nodes simulate accessibility.  The diagonal lemma supplies
the self-reference, and incompleteness prevents the base theory from
collapsing the simulated alternatives.  This paragraph records the
architecture, not the full arithmetization; the complete proof remains carried
by \cite{Solovay1976,Boolos1993}.
\end{proofboundary}

\begin{hbdefinition}{def:p05-rosser-provability}{Rosser provability and witness comparison}
At the metalinguistic level, let
\(\operatorname{ProofRel}_T(p,x)\) be the primitive-recursive relation saying
that \(p\) codes a \(T\)-proof with conclusion code \(x\), and let
\(\operatorname{NegCode}(x)\) be the primitive-recursive code of the negation
of the sentence coded by \(x\).  The first relation is the one represented by
the exact formula \(\mathsf{Proof}_T(p,x)\); the second is a function on
codes, not an undeclared term of \(\Lang_A\).

Form the primitive-recursive relation
\[
 \operatorname{RProofRel}_T(p,x)
 \quad\Longleftrightarrow\quad
 \operatorname{ProofRel}_T(p,x)\land
 \forall q\leq p\,
   \neg\operatorname{ProofRel}_T
      (q,\operatorname{NegCode}(x)).
\]
Bounded search preserves primitive recursiveness.  By the representability
theorem, fix a \(\Sigma^0_1\) formula
\(\mathsf{RProof}_T(p,x)\) numeralwise representing this relation in
\(\mathsf Q\), and define the object-language formula
\[
 \mathsf{RProv}_T(x)\defeq\exists p\,\mathsf{RProof}_T(p,x).
\]
Thus a Rosser proof of the sentence coded by \(x\) is strictly smaller than
every proof of its negation.  The metalinguistic names above are necessary to
construct the bounded comparison before choosing a representing formula; they
stand as a primitive-recursive relation and function on codes, differ from
their object-language representatives, and are used only in this construction.

More generally, a strict witness comparison says that some proof of \(A^*\)
has code \(p\) and no proof of \(B^*\) has code \(q\leq p\); its non-strict
variant uses \(q<p\).  Applying the same relation-first construction gives
\(\Sigma^0_1\) representatives of both comparisons.  No unqualified
\(\prec\) or \(\preccurlyeq\) glyph is assigned a new local role, and these
proof-code comparisons are not a second Kripke accessibility relation.  The
distinction is needed because ordinary \(\mathsf{Prov}_T\) forgets proof
order, and it is used by the comparison in Theorem~\ref{thm:rosser}.

Guaspari and Solovay's source names the resulting witness-comparison extension
of \(\mathsf{GL}\) the logic \(R\) and proves an arithmetical completeness
theorem for it \cite{GuaspariSolovay1979}.  This sentence is a bibliographic
pointer: the handbook does not introduce \(R\) as a local calculus or conflate
it with \(\mathsf{GL}\); see also \cite{Lindstrom2003}.
\index{Rosser provability}\index{witness comparison}
\end{hbdefinition}

\begingroup\small
\begin{longtable}{@{}p{0.26\textwidth}p{0.31\textwidth}p{0.33\textwidth}@{}}
\caption{Provability-logic dictionary.  The middle column is
the exact arithmetical reading; no row identifies Kripke accessibility with a
proof relation.}
\label{tab:p05-provability-dictionary}\\
\toprule
Modal principle & Arithmetical reading under \(^*\) & Governing result \\
\midrule
\endfirsthead
\toprule
Modal principle & Arithmetical reading under \(^*\) & Governing result \\
\midrule
\endhead
\(\Box(A\to B)\to(\Box A\to\Box B)\) &
\(\mathsf{Prv}^{19}_{\PA}(\overline{\ulcorner A^*\to B^*\urcorner})\to
(\mathsf{Prv}^{19}_{\PA}(\overline{\ulcorner A^*\urcorner})\to
 \mathsf{Prv}^{19}_{\PA}(\overline{\ulcorner B^*\urcorner}))\) &
D2, Theorem~\ref{thm:p05-direct-pa-internal-mp}. \\
\addlinespace
Theorem-only necessitation & If \(\PA\derives A^*\), then
\(\PA\derives\mathsf{Prv}^{19}_{\PA}
(\overline{\ulcorner A^*\urcorner})\). &
D1, Theorem~\ref{thm:p05-direct-pa-internal-mp}; theorem premises only. \\
\addlinespace
\(\Box A\to\Box\Box A\) &
\(\mathsf{Prv}^{19}_{\PA}(\overline{\ulcorner A^*\urcorner})\to
\mathsf{Prv}^{19}_{\PA}(\overline{\ulcorner
\mathsf{Prv}^{19}_{\PA}(\overline{\ulcorner A^*\urcorner})\urcorner})\) &
D3, Theorem~\ref{thm:p05-direct-pa-d3}. \\
\addlinespace
\(\Box(\Box A\to A)\to\Box A\) & Formalized local reflection for
\(A^*\) implies provability of \(A^*\). &
L\"ob's theorem. \\
\addlinespace
Modal fixed point & A formula modalized in \(p\) has a fixed point in
\(\mathsf{GL}\). & Modal counterpart of diagonal self-reference; the two
theorems are not literally the same construction. \\
\addlinespace
\(\mathsf{GL}\vdash A\) & \(\PA\vdash A^*\) for every fixed-\(\PA\)
arithmetical realization. & Theorem~\ref{thm:p05-solovay-completeness}. \\
\bottomrule
\end{longtable}
\endgroup
